\begin{frame}{Konstrukcija pravokotnice}
	\begin{columns}
		\begin{column}{0.55\textwidth}
			\begin{itemize}
				\onslide<1-7>{\item Dana sta premica \(p\) in točka \(T\).}
				\onslide<2-7>{\item Nariši lok \(k\) s središčem v \(T\).}
				\onslide<3-7>{\item Lok \(k\) seče premico \(p\) v točkah \(A\) in \(B\).}
				\onslide<4-7>{\item Nariši lok \(m\) s središčem v \(A\).}
				\onslide<5-7>{\item Nariši lok \(n\) s središčem v \(B\) in z enakim polmerom.}
				\onslide<6-7>{\item Loka se sečeta v točki \(C\).}
				\onslide<7-7>{\item Premica \(TC\) je pravokotna na \(p\).}
			\end{itemize}
		\end{column}
		\begin{column}{0.45\textwidth}
			\centering
			\includegraphics<1>[width=50mm]{slike/fig-1.png}%
			\includegraphics<2>[width=50mm]{slike/fig-2.png}%
			\includegraphics<3>[width=50mm]{slike/fig-3.png}%
			\includegraphics<4>[width=50mm]{slike/fig-4.png}%
			\includegraphics<5>[width=50mm]{slike/fig-5.png}%
			\includegraphics<6>[width=50mm]{slike/fig-6.png}%
			\includegraphics<7>[width=50mm]{slike/fig-7.png}%
		\end{column}
	\end{columns}
\end{frame}
